\section{Results}\label{sec:results}

\begin{guide}
\textbf{About 3--4 pages, including 4 tables and 2 figures.} Report what happened, in hypothesis order, and keep the ``why'' for \cref{sec:discussion}. The table numbers come from \texttt{Code/results/*.md} and \texttt{HANDOFF.md}; results are the mean $\pm$ std over 5 seeds. Give one paragraph per table or figure: what it shows, the key number, and which hypothesis it bears on.
\end{guide}

\subsection{Forecasting (H1, H2)}\label{sec:res-forecast}

\begin{table}[ht]
  \centering
  \caption{Next-day BTC log-return forecasts, 2025 test (mean $\pm$ std over 5 seeds). Naive ``no change'' MAPE is 1.785\%.}
  \label{tab:forecast}
  \small
  \begin{tabular}{llcc}
    \toprule
    Model & Inputs & RMSE & MAPE (\%) \\
    \midrule
    Vanilla LSTM            & univariate   & 0.02337 $\pm$ 0.00005 & 1.776 \\
    Vanilla GRU             & univariate   & 0.02348 $\pm$ 0.00019 & 1.783 \\
    Bi-LSTM                 & univariate   & 0.03355 $\pm$ 0.00983 & -- \\
    Bi-LSTM                 & multivariate & 0.02689 $\pm$ 0.00445 & 2.104 \\
    Alpha Agent (attention) & multivariate & 0.02372 $\pm$ 0.00051 & 1.816 \\
    \bottomrule
  \end{tabular}
\end{table}

\begin{guide}
One paragraph covering three points:
\begin{itemize}
  \item every model sits at about the naive baseline;
  \item \textbf{H1}: same Bi-LSTM, 19 features vs BTC alone; DM tests significant on 4 of 5 seeds but split 2--2 in direction, so it is supported on average but not consistently;
  \item \textbf{H2}: attention was significantly better on 2 of 5 seeds, both where the plain Bi-LSTM trained unstably, and a small non-significant edge otherwise.
\end{itemize}
Then one or two sentences on direction: the pooled hit rate was 50.2\% [46\%, 54\%] across 2024--2026, and no simple rule (always up, momentum, mean reversion) did better than a coin flip either. \texttt{phase5\_direction\_baselines.png} can go in the appendix.
\end{guide}
\tbc{}

\subsection{Volatility forecasting}\label{sec:res-vol}

\begin{figure}[ht]
  \centering
  \includegraphics[width=0.72\textwidth]{phase3_volatility_target.png}
  \caption{Skill relative to naive persistence for each volatility-target variant.}
  \label{fig:vol-target}
\end{figure}

\begin{guide}
One paragraph covering three points:
\begin{itemize}
  \item every 14-day variant lost heavily to persistence (\cref{fig:vol-target});
  \item with the single-day target, the Beta Agent's RMSE was 0.01974 against 0.02079 for persistence and 0.02246 for GARCH in 2025, and 0.0204 against 0.0250 in 2026;
  \item an honest caveat: the forecast is biased high (average 3.0\% against an actual 1.78\%) and correlates only 0.19 with the actual move.
\end{itemize}
Add one sentence on the jump classifier \tbc{calibration result from Phase 3}.
\end{guide}
\tbc{}

\subsection{Trading performance (H3)}\label{sec:res-h3}

\begin{table}[ht]
  \centering
  \caption{Backtest from R1{,}000 at a 2x leverage cap (mean over 5 runs). Bottom rows: Risk Engine minus fixed rules, paired block bootstrap, 95\% interval and one-sided $p$.}
  \label{tab:trading}
  \small
  \begin{tabular}{lcccccc}
    \toprule
    & \multicolumn{3}{c}{2025 test} & \multicolumn{3}{c}{2026 forward (unseen)} \\
    \cmidrule(lr){2-4}\cmidrule(lr){5-7}
    & Risk Engine & Fixed & BTC hold & Risk Engine & Fixed & BTC hold \\
    \midrule
    Final balance & R1{,}045 & R931 & R1{,}000 & R1{,}024 & R986 & R942 \\
    Lowest balance & R968 & R884 & R872 & R943 & R835 & R660 \\
    Max drawdown & 4.6\% & 17.2\% & 31.8\% & 8.7\% & 22.8\% & 39.6\% \\
    Sharpe & 0.74 & $-0.50$ & 0.19 & 0.35 & $-0.18$ & 0.04 \\
    Liquidations & 0 & 0 & -- & 0 & 0 & -- \\
    \midrule
    Drawdown cut (pp) & \multicolumn{3}{c}{+12.6 [7.0, 20.9], $p<0.001$} & \multicolumn{3}{c}{+14.1 [4.5, 19.2], $p<0.001$} \\
    Sharpe gain & \multicolumn{3}{c}{+1.24 [0.20, 2.45], $p=0.011$} & \multicolumn{3}{c}{+0.53 [$-0.44$, 1.65], $p=0.14$} \\
    \bottomrule
  \end{tabular}
\end{table}

\begin{figure}[ht]
  \centering
  \begin{subfigure}{0.49\textwidth}
    \includegraphics[width=\textwidth]{explain_equity_2025.png}
    \caption{2025 test}
  \end{subfigure}\hfill
  \begin{subfigure}{0.49\textwidth}
    \includegraphics[width=\textwidth]{explain_equity_2026.png}
    \caption{2026 forward}
  \end{subfigure}
  \caption{Account value from R1{,}000: Risk Engine, fixed rules, and buy-and-hold BTC. Lines are the average of 5 runs; shading spans best to worst run.}
  \label{fig:equity}
\end{figure}

\begin{guide}
Two paragraphs. The \textbf{first} is the headline: drawdown was cut by about three-quarters in both periods, and the cut was significant both times. The 2026 result is the cleanest evidence, because the data did not exist when the system was built. The Sharpe gain was significant in 2025 but not in 2026. The Risk Engine traded on 63 of 201 days in 2025, against 201 for fixed rules.

The \textbf{second} is robustness:
\begin{itemize}
  \item \textbf{Rule ablation} (\cref{tab:ablation}): no single rule carries the result, and the volatility stop matters most.
  \item \textbf{Leverage stress test:} max drawdown for the Risk Engine vs fixed rules was 12.6\% vs 38.3\% at 5x, 24.2\% vs 42.2\% at 10x, and 38.2\% vs 68.8\% at 20x. The cut is significant at every cap. However, the Risk Engine was liquidated twice at 20x and fixed rules never were, so the liquidation part of H3 is not supported.
  \item \textbf{Forecaster swap:} replacing the Alpha Agent with the Bi-LSTM or a univariate LSTM made no significant difference.
\end{itemize}
\end{guide}
\tbc{}

\begin{table}[ht]
  \centering
  \caption{Rule ablation, 2025 test: each rule replaced by its fixed version and the system retuned.}
  \label{tab:ablation}
  \small
  \begin{tabular}{lccc}
    \toprule
    Version & Max drawdown & Extra drawdown (pp) & Sharpe \\
    \midrule
    All five rules            & 4.6\%  & --    & 0.74 \\
    Fixed stop-loss           & 8.2\%  & +3.6  & $-0.14$ \\
    No weak-prediction filter & 7.9\%  & +3.3  & 0.39 \\
    Fixed position size       & 6.8\%  & +2.2  & 0.71 \\
    No jump veto              & 5.6\%  & +1.0  & 0.47 \\
    Fixed leverage            & 4.9\%  & +0.3  & 0.74 \\
    All fixed                 & 17.2\% & +12.6 & $-0.50$ \\
    \bottomrule
  \end{tabular}
\end{table}

\subsection{Summary}
\begin{table}[ht]
  \centering
  \caption{Outcome for each hypothesis.}
  \label{tab:scorecard}
  \small
  \begin{tabularx}{\textwidth}{lX}
    \toprule
    H1 & Partly supported: lower RMSE on average, not consistent across seeds. \\
    H2 & Partly supported: attention helps only when the plain Bi-LSTM trains unstably; no measurable effect in trading. \\
    H3 & Drawdown: supported in 2025 and 2026 and at every leverage cap. Risk-adjusted return: not worse (significant gain in 2025, not in 2026). Liquidations: not supported. \\
    \bottomrule
  \end{tabularx}
\end{table}
